\chapter{Introdução}
\section{Considerações Iniciais}
\par{A grafita, um mineral de relevância crescente no panorama tecnológico e industrial, destaca-se por suas propriedades únicas e aplicações versáteis, desde o desenvolvimento de nanomateriais, como o grafeno, até seu uso em produtos industriais diversos. Este projeto de pesquisa foca na exploração de jazidas de grafita utilizando técnicas avançadas de aprendizado de máquina e sensoriamento remoto, concentrando-se no sistema de nappes de Socorro-Guaxupé.}

\par{Exploramos como os métodos geofísicos, em conjunto com o processamento de dados aerogeofísicos, podem ser utilizados para identificar padrões associados à mineralização de grafita. O objetivo é desenvolver um modelo preditivo robusto que melhore a eficiência da prospecção mineral e a gestão dos recursos naturais.}

\par{Esta pesquisa se insere em um contexto onde a demanda por grafita está em ascensão devido às suas aplicações em tecnologias emergentes. A grafita é amplamente distribuída em diversos tipos de rochas, mas a identificação de depósitos economicamente viáveis é um desafio que requer abordagens inovadoras. A integração de dados geofísicos com técnicas de aprendizado de máquina apresenta uma oportunidade única para abordar essa questão, potencializando a identificação de áreas promissoras para exploração mineral.}

\section{Localização e Acessos}
\par{A região de Socorro está situada no nordeste do estado de São Paulo, próxima à divisa com Minas Gerais, a uma distância aproximada de 130 km da capital paulista. O acesso à cidade pode ser feito principalmente por meio das rodovias Anhanguera (SP-330) e Bandeirantes (SP-348), que ligam São Paulo à região de Jundiaí, onde se acessa a Rodovia João Cereser e, em seguida, a Rodovia Engenheiro Constâncio Cintra. De lá, o trajeto continua pelas rodovias Perimetral de Itatiba e Alkindar Monteiro Junqueira até Bragança Paulista, onde se conecta com a Rodovia Capitão Barduíno (BR-146), que leva diretamente a Socorro.}

\par{Este trajeto inclui algumas rodovias com pedágios e leva aproximadamente 2 horas e 40 minutos, dependendo das condições de tráfego. O mapa fornecido ilustra o percurso, destacando as principais vias de acesso entre o Instituto de Geociências da USP e a cidade de Socorro.}

\section{Contexto Geológico}

\par{A região de Socorro insere-se em um contexto geológico complexo, compondo, juntamente com o Complexo Guaxupé, uma unidade paleoproterozóica a arqueana denominada Placa São Paulo, ou Nappe Socorro-Guaxupé, de acordo com \cite{campos_neto_2004}. Este terreno foi intensamente retrabalhado durante o Neoproterozóico, quando processos tectônicos de colisão e subducção influenciaram sua configuração atual \cite{silva_2005}.}

As rochas presentes na área incluem granulitos, gnaisses e granitos, que foram afetados pela Zona de Cisalhamento de Socorro. Essa zona representa um importante limite tectônico e foi responsável por contatos tectônicos entre rochas de alto grau metamórfico e unidades graníticas ao leste e migmatíticas ao oeste \cite{oliveira_2010}.

A evolução tectônica dessa área é marcada por uma trajetória metamórfica em alta pressão, seguida de uma descompressão quase isotermal, evidenciada nas paragêneses minerais das rochas granulíticas. As condições de pressão e temperatura (P-T) registradas nas assembléias de granada e ortopiroxênio sugerem um contexto de alto fluxo térmico, possivelmente relacionado ao magmatismo granítico da região \cite{souza_2012, juliani_2014}.

A Zona de Cisalhamento de Socorro também apresenta características transpressivas, com componente de cavalgamento e deformações que variaram de rúptil-dúctil a dúctil, provocando intercalações litológicas complexas e reequilíbrios metamórficos nas rochas da região. Esses processos tectônicos foram fundamentais para a exumação dos granulitos e para a sua estruturação atual, evidenciando o caráter policíclico do terreno estudado \cite{freitas_2018}.
